\section{Estructura de Archivos}

\subsection{Regla: una clase de nivel superior por archivo}

\textbf{Regla:} Cada archivo contiene exactamente una clase, interfaz, \texttt{record} o \texttt{enum} de nivel superior; el nombre del archivo coincide con el del tipo (McConnell, 2004).

\textbf{Código correcto (archivo \texttt{TurnManager.cs}):}
\begin{lstlisting}[style=csharpstyle]
namespace BoardGame.Logic;

public sealed class TurnManager
{
}
\end{lstlisting}

\textbf{Código incorrecto (archivo \texttt{GameLogic.cs}):}
\begin{lstlisting}[style=csharpstyle]
namespace BoardGame.Logic;

public sealed class TurnManager
{
}

public sealed class ScoreBoard
{
}
\end{lstlisting}

\subsection{Regla: orden de los elementos dentro de la clase}

\textbf{Regla:} Orden fijo dentro de una clase: constantes, campos estáticos, campos de instancia, constructores, propiedades, métodos (McConnell, 2004).

\textbf{Código correcto:}
\begin{lstlisting}[style=csharpstyle]
public sealed class DiceRoller
{
    private const int MinimumFaceValue = 1;

    private static readonly Random _random = Random.Shared;

    private readonly int _maximumFaceValue;

    public DiceRoller(int maximumFaceValue)
    {
        _maximumFaceValue = maximumFaceValue;
    }

    public int LastRoll { get; private set; }

    public int Roll()
    {
        LastRoll = _random.Next(MinimumFaceValue, _maximumFaceValue + 1);
        return LastRoll;
    }
}
\end{lstlisting}

\textbf{Código incorrecto:}
\begin{lstlisting}[style=csharpstyle]
public sealed class DiceRoller
{
    public int Roll()
    {
        LastRoll = _random.Next(MinimumFaceValue, _maximumFaceValue + 1);
        return LastRoll;
    }

    private readonly int _maximumFaceValue;

    public int LastRoll { get; private set; }

    public DiceRoller(int maximumFaceValue)
    {
        _maximumFaceValue = maximumFaceValue;
    }

    private const int MinimumFaceValue = 1;
    private static readonly Random _random = Random.Shared;
}
\end{lstlisting}

\subsection{Regla: orden y continuidad de directivas \texttt{using}}

\textbf{Regla:} Las directivas \texttt{using} se colocan fuera del espacio de nombres en un único bloque sin líneas en blanco entre grupos. El bloque se ordena en tres grupos consecutivos, cada uno alfabético internamente: primero los espacios de nombres de \texttt{System}, después las bibliotecas de terceros (por ejemplo \texttt{log4net}) y al final los espacios de nombres propios del proyecto (\texttt{BoardGame.*}) (Microsoft, 2022a; Microsoft, 2022b). La configuración equivalente es:

\begin{itemize}
  \item \texttt{csharp\_using\_directive\_placement = outside\_namespace}
  \item \texttt{dotnet\_sort\_system\_directives\_first = true}
  \item \texttt{dotnet\_separate\_import\_directive\_groups = false}
\end{itemize}

\textbf{Código correcto:}
\begin{lstlisting}[style=csharpstyle]
using System;
using System.Collections.Generic;
using log4net;
using BoardGame.Core;

namespace BoardGame.Services;
\end{lstlisting}

\textbf{Código incorrecto:}
\begin{lstlisting}[style=csharpstyle]
using log4net;

using System.Collections.Generic;
using BoardGame.Core;
using System;

namespace BoardGame.Services;
\end{lstlisting}

\subsection{Regla: agrupación mediante espacios de nombres}

\textbf{Regla:} Todo tipo pertenece a un espacio de nombres que refleja la organización del proyecto y la responsabilidad del archivo. Se usa la declaración de ámbito de archivo y un solo espacio de nombres por archivo; por ejemplo, los servicios de guardado se agrupan bajo \texttt{BoardGame.Services.Saves} (Microsoft, 2026d).

\textbf{Código correcto (archivo \texttt{SaveGameService.cs}):}
\begin{lstlisting}[style=csharpstyle]
namespace BoardGame.Services.Saves;

public sealed class SaveGameService
{
}
\end{lstlisting}

\textbf{Código incorrecto (archivo \texttt{SaveGameService.cs}):}
\begin{lstlisting}[style=csharpstyle]
namespace BoardGame;

public sealed class SaveGameService
{
}
\end{lstlisting}
